\addchap{\lsPrefaceTitle}
For almost 150 years a big question mark has hung over how the Austronesian languages of what is now eastern Indonesia
and Timor-Leste fit into the wider Austronesian world,
and how they got there.
The area is among the most severely under-documented regions of Austronesia.
There has been no consensus on the classification of these languages,
and scholars looking at selected issues in parts of the region admit to being perplexed at how they all fit together.
This study is the first systematic attempt at subgrouping the whole region based primarily on the historical phonology,
supplemented by issues in the morphosyntax
and the lexicon.
The study also benefits from insights into the prehistory of the region from archaeology
and DNA studies,
and from an awareness of the significance of intense long-term contact with Papuan languages.

Based on evidence of shared innovations in their historical phonology,
this study identifies nine subgroups in Linguistic Wallacea (one of which links to Celebic in Sulawesi),
along with the internal structure of these subgroups (chapter \ref{ch:BL}--\ref{ch:Ban}).
In doing so,
this study also sheds light on the classification of a number of languages whose classification has been previously unclear.
The discontinuities in the historical phonology suggest that different groups
of people speaking different Austronesian languages got off different boats at different places,
and probably at different times.
The study finds no evidence in the historical phonology
and morphosyntax to support a monolithic Austronesian advance-by-wave through the region,
nor evidence of a common Austronesian parent language below PMP that link all Wallacean subgroups together.

While it has long been known that the Austronesian languages of this region are somehow “different”,
only recently has the significance of the intense contact over long periods
of time with non-Austronesian Papuan languages become increasingly recognised.
Speakers of typologically SVO Austronesian languages with prepositions,
preverbal negation, numbers before nouns,
and post-posed possessors came into contact in the region with speakers
of different languages of several unrelated Papuan families,
with postpositions, clause-final negation,
numbers following nouns,
preposed possessors, and other features typically associated with SOV languages.

The Austronesian languages in the region accommodated to adopt these features but did not do so in a uniform way,
such that the features attributed to contact are distributed unevenly across the region.
Some are not found in some areas or subgroups,
and their presence varies within subgroups.
Several distribution maps of phonological,
grammatical, and lexical features show that many features are not found in all subgroups,
do not align with each other,
and some are found outside the region.
The Austronesian languages in the target region are a kind of uneven hybrid that make them grammatically
and typologically different from Austronesian languages to the west
and north.

The study evaluates several issues that have been proposed over the years,
along with some new possibilities to link languages
and subgroups in the region together in different ways,
but finds no features or sound changes within the region that qualify as exclusively
shared innovations inherited from a common parent language.
Possible scenarios are explored of how Austronesian-speaking peoples came into the region.
The methodological problem of explaining why typological features
and lexical items are distributed unevenly in the region is also addressed.
The implications of this study are many,
and include a revised picture of the Austronesian-speaking world.

Several factors combined to enable this more in-depth study
of languages in the region than has been previously possible.
Both authors have extensive experience in the region.
Many Dutch-era sources have recently become accessible online.
And several recent publications
and unpublished data have been generously shared by other scholars.
All this enabled the authors to glean data from at l\ili{east} 520 Austronesian
and Papuan languages/speech varieties from within the region as well as to the west
and \ili{east} of it,
providing context.
Within the target region itself,
data have been gleaned from at l\ili{east} 296 varieties (260 Austronesian, 36 Papuan),
some of which are now extinct.
The volume is data rich with 317 data tables,
79 figures (including 32 maps),
and 235 numbered examples/lists.
 
 